\begin{table}[!h]
\centering
\caption{\label{tab:recap_holobiotique}Effectifs annuels des poissons holobiotiques comptabilisés sur les stations de contrôle des migrations du bassin de la Seine}
\centering
\resizebox{\ifdim\width>\linewidth\linewidth\else\width\fi}{!}{
\begin{tabular}[t]{>{\raggedright\arraybackslash}p{4cm}C{2cm}C{2cm}C{2cm}}
\toprule
Espèces & Poses & Pontoise & Carandeau\\
\midrule
\cellcolor{gray!10}{Ablette} & \cellcolor{gray!10}{134} & \cellcolor{gray!10}{15 005} & \cellcolor{gray!10}{12 154}\\
Amour blanc & 0 & 8 & 3\\
\cellcolor{gray!10}{Aspe} & \cellcolor{gray!10}{2} & \cellcolor{gray!10}{2} & \cellcolor{gray!10}{}\\
Barbeau fluviatile & 138 & 1 916 & 266\\
\cellcolor{gray!10}{Brochet} & \cellcolor{gray!10}{11} & \cellcolor{gray!10}{28} & \cellcolor{gray!10}{8}\\
Brème commune & 5 320 & 1 682 & 391\\
\cellcolor{gray!10}{Carassin commun} & \cellcolor{gray!10}{} & \cellcolor{gray!10}{} & \cellcolor{gray!10}{5}\\
Carpe commune & 26 &  & 22\\
\cellcolor{gray!10}{Chevesne} & \cellcolor{gray!10}{9} & \cellcolor{gray!10}{667} & \cellcolor{gray!10}{836}\\
Crabe chinois & 220 &  & \\
\cellcolor{gray!10}{Cyprinidé indéterminé} & \cellcolor{gray!10}{66 062} & \cellcolor{gray!10}{3 433} & \cellcolor{gray!10}{3 014}\\
Gardon &  & 4 431 & 35\\
\cellcolor{gray!10}{Gobie à tâches noires} & \cellcolor{gray!10}{6} & \cellcolor{gray!10}{} & \cellcolor{gray!10}{}\\
Hotu & 16 & 9 & 20\\
\cellcolor{gray!10}{Ide mélanote} & \cellcolor{gray!10}{1 372} & \cellcolor{gray!10}{} & \cellcolor{gray!10}{1}\\
Lamproie indéterminée &  & 26 & \\
\cellcolor{gray!10}{Perche} & \cellcolor{gray!10}{58} & \cellcolor{gray!10}{208} & \cellcolor{gray!10}{20}\\
Sandre & 1 &  & 1\\
\cellcolor{gray!10}{Silure glane} & \cellcolor{gray!10}{267} & \cellcolor{gray!10}{205} & \cellcolor{gray!10}{74}\\
Tanche & 5 & 1 & 4\\
\cellcolor{gray!10}{Taxon indéterminé} & \cellcolor{gray!10}{} & \cellcolor{gray!10}{310} & \cellcolor{gray!10}{}\\
Truite fario & 2 &  & \\
\cellcolor{gray!10}{Vandoise} & \cellcolor{gray!10}{} & \cellcolor{gray!10}{23} & \cellcolor{gray!10}{40}\\
\bottomrule
\end{tabular}}
\end{table}
